\documentclass[10pt,a4paper]{amsart}
\usepackage[margin=19mm]{geometry}
\usepackage{amsmath,amssymb,mathtools}
\usepackage[scr=rsfs,cal=euler]{mathalfa}
\usepackage{xfrac}
\usepackage[unicode,hidelinks,bookmarks=false]{hyperref}
\newcommand{\ie}{i.e.\ }
\newcommand{\cf}{cf.\ }
\newcommand{\resp}{respectively\ }
\newcommand{\Subsection}{\subsection}
\providecommand{\nobreakdash}{\nobreak\hbox{-}}
\setlength{\emergencystretch}{2em}
\multlinegap0pt
\usepackage{xcolor}
\usepackage{amssymb}
\usepackage[shortlabels]{enumitem}
\usepackage{tikz}
\usepackage{float}
\usepackage{mathtools}
\usepackage{dsfont}
\newtheorem{lemma}{Lemma}
\title{behaviour of symmetric power $L$-function of a Holomorphic cusp form over certain sum of squares}
\author{Jewel Mahajan, Arnab Mitra}
\date{Source: arXiv:2606.30603v1 (CC BY 4.0)}
\begin{document}
\maketitle

\noindent\emph{Source and licence.} behaviour of symmetric power $L$-function of a Holomorphic cusp form over certain sum of squares. Version arXiv:2606.30603v1 (CC BY 4.0), distributed by arXiv under the Creative Commons Attribution 4.0 International licence, http://creativecommons.org/licenses/by/4.0/, as recorded in the arXiv licence metadata for this exact version and shown on the abstract page https://arxiv.org/abs/2606.30603v1. Full licence notice: \texttt{licenses/arxiv-2606.30603-CC-BY-4.0.txt}.\par\smallskip

\noindent\textit{Editorial scope.} Counted unit: the article's lemma lambdasquare (source0.tex lines 390--395). The article's complete original proof of it is delivered with it (source0.tex lines 396--421, one \texttt{proof} environment of 26 lines). No mathematics is written, repaired or re-derived: the statement and the proof are the article's own lines, in the article's own notation, and no retained line is substituted or re-flowed.  No further source-local context is needed: every cross-reference of every retained line resolves inside the two delivered spans.  Nothing was omitted from any retained span, so the package carries no omission record.  No retained line contains a citation, so the wrapper carries no bibliography and no bibliography entry.  No retained line contains a figure environment, an image-inclusion command or drawing code, so no image is delivered and the package declares no asset dependency.  Editorial formatting only, in addition: an AMS \texttt{amsart} preamble that restates the article's own declarations and macros used by the retained lines, namely the environments \texttt{lemma} on separate counters, together with the standard packages those lines need.  The retained environments print in this document as Lemma 1 in order of appearance, whereas the article prints the same items with the numbering of its own full document.  The complete native source of the article as distributed in the arXiv e-print for this exact version is bundled unchanged as \texttt{sources/204038/source0.tex}.
\begin{lemma}\label{lambdasquare}
Let $f$ be a Hecke eigenform with Satake parameters $\alpha_p,\beta_p$ satisfying $\alpha_p\beta_p=1$.
For any integer $j \ge 1$ and any prime $p$,
\[
\lambda_{\mathrm{sym}^{j}f}^{2}(p) \;=\; 1 \;+\; \sum_{\ell=1}^{j} \lambda_{\mathrm{sym}^{2\ell}f}(p).\]
\end{lemma}

\begin{proof}
Write $\lambda := \lambda_{\mathrm{sym}^{j}f}(p) = \sum_{m=0}^{j} \alpha_p^{\,j-m}\beta_p^{\,m}$.
Then, using $\beta_p = \alpha_p^{-1}$, we obtain
\[
\lambda^{2} = \sum_{m=0}^{j}\sum_{m'=0}^{j} \alpha_p^{\,2j-(m+m')}\beta_p^{\,m+m'}
      = \sum_{t=0}^{2j} N_t \, \alpha_p^{\,2j-t}\beta_p^{\,t}= \sum_{t=0}^{2j} N_t \, \alpha_p^{\,2j-2t},
\]
where $N_t$ counts pairs $(m,m')$ with $m+m'=t$ and $0\le m,m'\le j$.
One has $N_t = t+1$ for $0\le t\le j$ and $N_t = 2j-t+1$ for $j<t\le 2j$. In particular, $N_t = N_{2j-t}$ for $0\le t\le 2j$. Therefore, 
\begin{align}
    \lambda^{2} = &\sum_{t=0}^{2j} N_t \, \alpha_p^{\,2j-2t}\\
      =&N_{j}+ \sum_{t=0}^{j-1} N_t \, \alpha_p^{\,2j-2t}+ \sum_{t=j+1}^{2j} N_{2j-t} \, \alpha_p^{\,2j-2t}\quad (\text{since } N_t = N_{2j-t})\\
      =&N_{j}+ \sum_{t=0}^{j-1} N_t \, \alpha_p^{\,2j-2t}+ \sum_{t=0}^{j-1} N_{t} \, \alpha_p^{\,-2j+2t}\\
      =&N_{j}+ \sum_{k=1}^{j} N_{j-k} \, \left(\alpha_p^{\,2k}+\alpha_p^{\,-2k}\right)\\
      =& (j+1)+ \sum_{k=1}^{j} (j-k+1) \, \left(\alpha_p^{\,2k}+\alpha_p^{\,-2k}\right).
\end{align}
On the other hand, we have
\[
\sum_{l=1}^{j}\lambda_{\mathrm{sym}^{2\ell}f}(p) = \sum_{l=1}^{j}\sum_{u=0}^{2\ell} \alpha_p^{\,2\ell-u}\beta_p^{\,u}= \sum_{l=1}^{j}\sum_{u=0}^{2\ell} \alpha_p^{\,2\ell-2u} = \sum_{l=1}^{j}\sum_{m=-\ell}^{\ell} \alpha_p^{\,2m}.
\]
In the double sum, for \( k \geq 1 \), the term \( \alpha_p^{\,2k} \) (and similarly \( \alpha_p^{\,-2k} \)) occurs once for each \( \ell \) satisfying \( \ell \ge k \), that is, for \( \ell = k, k+1, \dots, j \). Thus, the coefficient of \( \alpha_p^{\,2k} \) or \( \alpha_p^{\,-2k} \) is \( j-k+1 \).
The constant term $\alpha_p^{0}=1$ occurs in every $\lambda_{\operatorname{sym}^{2\ell} f}(p)$ for $\ell = 1, \dots, j$, giving total multiplicity $j$. Therefore,
\[
\sum_{l=1}^{j}\lambda_{\mathrm{sym}^{2\ell}f}(p) = \sum_{l=1}^{j}\sum_{m=-\ell}^{\ell} \alpha_p^{\,2m}=j+ \sum_{k=1}^{j} (j-k+1) \, \left(\alpha_p^{\,2k}+\alpha_p^{\,-2k}\right).
\]This completes the proof.
\end{proof}

\end{document}
